% =====================================================================
%  minggu-09.tex  --  Slide Pertemuan Minggu ke-9
%  Pengantar Kriptografi (MPFA014) -- Prodi Matematika FMIPA UBB
%
%  Topik silabus : Block cipher, DES, dan Triple DES
%  Sumber materi : buku_ajar/chapters/06 (Bab 6) dan chapters/07 (Bab 7)
%  Sub-CPMK      : 1.3
%
%  Cakupan: mode ECB, CBC, CFB, OFB, CTR; confusion dan diffusion;
%  jaringan Feistel; DES; Double DES; Triple DES. AES dibahas di Minggu 10.
%
%  Kompilasi dari dalam folder ppt/:
%      pdflatex -interaction=nonstopmode minggu-09.tex   (dua kali)
% =====================================================================

\documentclass[aspectratio=169,11pt]{beamer}
% =====================================================================
%  beamer-preamble.tex  --  prasarana dan gaya slide kuliah
%
%  Mata kuliah : Pengantar Kriptografi (MPFA014)
%  Program     : Program Studi Matematika, FMIPA, Universitas Bale Bandung
%
%  Berkas ini TIDAK memuat \documentclass. Setiap berkas minggu memuatnya
%  sendiri lalu menarik berkas ini:
%
%      \documentclass[aspectratio=169,11pt]{beamer}
%      % =====================================================================
%  beamer-preamble.tex  --  prasarana dan gaya slide kuliah
%
%  Mata kuliah : Pengantar Kriptografi (MPFA014)
%  Program     : Program Studi Matematika, FMIPA, Universitas Bale Bandung
%
%  Berkas ini TIDAK memuat \documentclass. Setiap berkas minggu memuatnya
%  sendiri lalu menarik berkas ini:
%
%      \documentclass[aspectratio=169,11pt]{beamer}
%      % =====================================================================
%  beamer-preamble.tex  --  prasarana dan gaya slide kuliah
%
%  Mata kuliah : Pengantar Kriptografi (MPFA014)
%  Program     : Program Studi Matematika, FMIPA, Universitas Bale Bandung
%
%  Berkas ini TIDAK memuat \documentclass. Setiap berkas minggu memuatnya
%  sendiri lalu menarik berkas ini:
%
%      \documentclass[aspectratio=169,11pt]{beamer}
%      \input{beamer-preamble}
%      \title{...} \subtitle{...} \setWeekInfo{Minggu n}{Topik}
%
%  Berkas minggu hanya berisi \setWeekInfo dan lingkungan frame.
%  Kompilasi dijalankan DARI DALAM folder ppt/ (lihat compile-ppt.bat),
%  karena gambar buku dirujuk dengan jalur relatif ../buku_ajar/.
% =====================================================================

% ---------------------------------------------------------------- bahasa
\usepackage[utf8]{inputenc}
\usepackage[T1]{fontenc}
\usepackage[indonesian]{babel}
\usepackage{lmodern}
\usepackage{microtype}

% ------------------------------------------------------------- matematika
\usepackage{amsmath,amssymb,mathtools}

% ----------------------------------------------------------------- tabel
\usepackage{booktabs}
\usepackage{array}
\usepackage{tabularx}

% --------------------------------------------------------------- grafik
\usepackage{tikz}
\usetikzlibrary{arrows.meta,positioning,shapes.geometric,calc,fit,backgrounds}

% ---------------------------------------------------------- kode program
\usepackage{listings}
\lstdefinestyle{slidecode}{
    basicstyle=\ttfamily\scriptsize,
    keywordstyle=\color{ubbblue}\bfseries,
    commentstyle=\color{ubbgreen}\itshape,
    stringstyle=\color{ubbred},
    numbers=left,
    numberstyle=\tiny\color{gray},
    numbersep=6pt,
    frame=single,
    rulecolor=\color{ubbblue!40},
    backgroundcolor=\color{ubbblue!3},
    breaklines=true,
    showstringspaces=false,
    tabsize=4,
    columns=fullflexible,
    keepspaces=true,
}
\lstset{style=slidecode}

% =====================================================================
%  WARNA INSTITUSI
% =====================================================================
\definecolor{ubbblue}{RGB}{12, 62, 120}      % warna utama
\definecolor{ubbteal}{RGB}{0, 118, 110}      % aksen kedua
\definecolor{ubbgreen}{RGB}{22, 122, 74}     % untuk blok contoh
\definecolor{ubbred}{RGB}{168, 42, 42}       % untuk blok peringatan
\definecolor{ubbgray}{RGB}{70, 76, 84}       % teks sekunder

% =====================================================================
%  IDENTITAS MATA KULIAH
% =====================================================================
\newcommand{\CourseName}{Pengantar Kriptografi}
\newcommand{\CourseCode}{MPFA014}
\newcommand{\CourseSemester}{Semester Ganjil}
\newcommand{\CourseProgram}{Program Studi Matematika}
\newcommand{\CourseFaculty}{Fakultas Matematika dan Ilmu Pengetahuan Alam}
\newcommand{\CourseInstitution}{Universitas Bale Bandung}
\newcommand{\CourseLecturer}{Cecep Suwanda, S.Si., M.Kom.}

% Label minggu, disetel tiap berkas lewat \setWeekInfo
\newcommand{\WeekLabel}{}
\newcommand{\WeekTopic}{}
\newcommand{\setWeekInfo}[2]{%
    \renewcommand{\WeekLabel}{#1}%
    \renewcommand{\WeekTopic}{#2}%
}

% =====================================================================
%  TEMA: header, footer, warna blok
% =====================================================================
\usetheme{default}
\setbeamertemplate{navigation symbols}{}
\setbeamertemplate{blocks}[rounded][shadow=false]

\setbeamercolor{palette primary}{bg=ubbblue, fg=white}
\setbeamercolor{palette secondary}{bg=ubbgray, fg=white}
\setbeamercolor{palette tertiary}{bg=ubbteal, fg=white}

\setbeamercolor{structure}{fg=ubbblue}
\setbeamercolor{normal text}{fg=black!92}
\setbeamercolor{alerted text}{fg=ubbred}

\setbeamercolor{block title}{bg=ubbblue, fg=white}
\setbeamercolor{block body}{bg=ubbblue!6, fg=black!92}
\setbeamercolor{block title alerted}{bg=ubbred, fg=white}
\setbeamercolor{block body alerted}{bg=ubbred!6, fg=black!92}
\setbeamercolor{block title example}{bg=ubbgreen, fg=white}
\setbeamercolor{block body example}{bg=ubbgreen!6, fg=black!92}

\setbeamercolor{itemize item}{fg=ubbblue}
\setbeamercolor{itemize subitem}{fg=ubbteal}
\setbeamercolor{enumerate item}{fg=ubbblue}

\setbeamerfont{headline}{size=\scriptsize}
\setbeamerfont{footline}{size=\scriptsize}
\setbeamerfont{block title}{size=\normalsize}

% Badan blok rata kiri (bukan rata kanan-kiri): teks slide lebih mudah dibaca
% dan baris terakhir yang pendek tidak lagi memicu "Underfull \hbox".
\addtobeamertemplate{block body begin}{\raggedright}{}
\addtobeamertemplate{block alerted body begin}{\raggedright}{}
\addtobeamertemplate{block example body begin}{\raggedright}{}

% Header: nama mata kuliah dan kode MK
\setbeamertemplate{headline}{%
    \begin{beamercolorbox}[wd=\paperwidth,ht=2.6ex,dp=1.2ex,leftskip=1em,rightskip=1em]{palette primary}%
        \usebeamerfont{headline}\CourseName\ \textperiodcentered\ \CourseCode%
        \hfill \CourseProgram%
    \end{beamercolorbox}%
}

% Footer: minggu dan nomor halaman
\setbeamertemplate{footline}{%
    \begin{beamercolorbox}[wd=\paperwidth,ht=2.4ex,dp=1.2ex,leftskip=1em,rightskip=1em]{palette secondary}%
        \usebeamerfont{footline}\WeekLabel\ \textperiodcentered\ \WeekTopic%
        \hfill \insertframenumber\,/\,\inserttotalframenumber%
    \end{beamercolorbox}%
}

% Halaman judul: judul, subjudul, lalu blok identitas mata kuliah
\setbeamertemplate{title page}{%
    \begin{centering}
        \vskip1.0em
        {\usebeamerfont{title}\usebeamercolor[fg]{title}\inserttitle\par}
        \vskip0.3em
        {\usebeamerfont{subtitle}\usebeamercolor[fg]{subtitle}\insertsubtitle\par}
        \vskip1.2em
        \textcolor{ubbblue}{\rule{0.72\textwidth}{0.8pt}}
        \vskip1.2em
        {\small
            \textbf{\CourseName}\ (\CourseCode)\par
            \vskip0.4em
            \CourseProgram\par
            \CourseFaculty\par
            \CourseInstitution\par
            \vskip0.6em
            \CourseSemester\ \textperiodcentered\ \WeekLabel\par
            \vskip0.8em
            Dosen Pengampu: \CourseLecturer\par
        }
        \vskip1.0em
        \textcolor{ubbblue}{\rule{0.72\textwidth}{0.4pt}}
    \end{centering}
}

% =====================================================================
%  PERINTAH BANTU
% =====================================================================
\newcommand{\R}{\mathbb{R}}
\newcommand{\N}{\mathbb{N}}
\newcommand{\Z}{\mathbb{Z}}
\newcommand{\code}[1]{\texttt{#1}}

% Rujukan gambar dari buku: \bookfig[lebar]{NN}{namaberkas.png}
% Jalur relatif terhadap folder ppt/ (dikompilasi dari sana).
\newcommand{\bookfig}[3][0.80\linewidth]{%
    \includegraphics[width=#1]{../buku_ajar/chapters/#2/figures/#3}%
}

% Keterangan sumber gambar, sepadan dengan \sumber di buku
\newcommand{\figsumber}[1]{%
    \par\vspace{1pt}{\tiny\itshape Sumber: #1\par}%
}

% Penanda blok materi yang sering dipakai di slide
\newcommand{\slidejudul}[1]{{\large\bfseries\textcolor{ubbblue}{#1}}\par\vspace{2pt}}
\newcommand{\kunci}[1]{\textcolor{ubbred}{\textbf{#1}}}

%      \title{...} \subtitle{...} \setWeekInfo{Minggu n}{Topik}
%
%  Berkas minggu hanya berisi \setWeekInfo dan lingkungan frame.
%  Kompilasi dijalankan DARI DALAM folder ppt/ (lihat compile-ppt.bat),
%  karena gambar buku dirujuk dengan jalur relatif ../buku_ajar/.
% =====================================================================

% ---------------------------------------------------------------- bahasa
\usepackage[utf8]{inputenc}
\usepackage[T1]{fontenc}
\usepackage[indonesian]{babel}
\usepackage{lmodern}
\usepackage{microtype}

% ------------------------------------------------------------- matematika
\usepackage{amsmath,amssymb,mathtools}

% ----------------------------------------------------------------- tabel
\usepackage{booktabs}
\usepackage{array}
\usepackage{tabularx}

% --------------------------------------------------------------- grafik
\usepackage{tikz}
\usetikzlibrary{arrows.meta,positioning,shapes.geometric,calc,fit,backgrounds}

% ---------------------------------------------------------- kode program
\usepackage{listings}
\lstdefinestyle{slidecode}{
    basicstyle=\ttfamily\scriptsize,
    keywordstyle=\color{ubbblue}\bfseries,
    commentstyle=\color{ubbgreen}\itshape,
    stringstyle=\color{ubbred},
    numbers=left,
    numberstyle=\tiny\color{gray},
    numbersep=6pt,
    frame=single,
    rulecolor=\color{ubbblue!40},
    backgroundcolor=\color{ubbblue!3},
    breaklines=true,
    showstringspaces=false,
    tabsize=4,
    columns=fullflexible,
    keepspaces=true,
}
\lstset{style=slidecode}

% =====================================================================
%  WARNA INSTITUSI
% =====================================================================
\definecolor{ubbblue}{RGB}{12, 62, 120}      % warna utama
\definecolor{ubbteal}{RGB}{0, 118, 110}      % aksen kedua
\definecolor{ubbgreen}{RGB}{22, 122, 74}     % untuk blok contoh
\definecolor{ubbred}{RGB}{168, 42, 42}       % untuk blok peringatan
\definecolor{ubbgray}{RGB}{70, 76, 84}       % teks sekunder

% =====================================================================
%  IDENTITAS MATA KULIAH
% =====================================================================
\newcommand{\CourseName}{Pengantar Kriptografi}
\newcommand{\CourseCode}{MPFA014}
\newcommand{\CourseSemester}{Semester Ganjil}
\newcommand{\CourseProgram}{Program Studi Matematika}
\newcommand{\CourseFaculty}{Fakultas Matematika dan Ilmu Pengetahuan Alam}
\newcommand{\CourseInstitution}{Universitas Bale Bandung}
\newcommand{\CourseLecturer}{Cecep Suwanda, S.Si., M.Kom.}

% Label minggu, disetel tiap berkas lewat \setWeekInfo
\newcommand{\WeekLabel}{}
\newcommand{\WeekTopic}{}
\newcommand{\setWeekInfo}[2]{%
    \renewcommand{\WeekLabel}{#1}%
    \renewcommand{\WeekTopic}{#2}%
}

% =====================================================================
%  TEMA: header, footer, warna blok
% =====================================================================
\usetheme{default}
\setbeamertemplate{navigation symbols}{}
\setbeamertemplate{blocks}[rounded][shadow=false]

\setbeamercolor{palette primary}{bg=ubbblue, fg=white}
\setbeamercolor{palette secondary}{bg=ubbgray, fg=white}
\setbeamercolor{palette tertiary}{bg=ubbteal, fg=white}

\setbeamercolor{structure}{fg=ubbblue}
\setbeamercolor{normal text}{fg=black!92}
\setbeamercolor{alerted text}{fg=ubbred}

\setbeamercolor{block title}{bg=ubbblue, fg=white}
\setbeamercolor{block body}{bg=ubbblue!6, fg=black!92}
\setbeamercolor{block title alerted}{bg=ubbred, fg=white}
\setbeamercolor{block body alerted}{bg=ubbred!6, fg=black!92}
\setbeamercolor{block title example}{bg=ubbgreen, fg=white}
\setbeamercolor{block body example}{bg=ubbgreen!6, fg=black!92}

\setbeamercolor{itemize item}{fg=ubbblue}
\setbeamercolor{itemize subitem}{fg=ubbteal}
\setbeamercolor{enumerate item}{fg=ubbblue}

\setbeamerfont{headline}{size=\scriptsize}
\setbeamerfont{footline}{size=\scriptsize}
\setbeamerfont{block title}{size=\normalsize}

% Badan blok rata kiri (bukan rata kanan-kiri): teks slide lebih mudah dibaca
% dan baris terakhir yang pendek tidak lagi memicu "Underfull \hbox".
\addtobeamertemplate{block body begin}{\raggedright}{}
\addtobeamertemplate{block alerted body begin}{\raggedright}{}
\addtobeamertemplate{block example body begin}{\raggedright}{}

% Header: nama mata kuliah dan kode MK
\setbeamertemplate{headline}{%
    \begin{beamercolorbox}[wd=\paperwidth,ht=2.6ex,dp=1.2ex,leftskip=1em,rightskip=1em]{palette primary}%
        \usebeamerfont{headline}\CourseName\ \textperiodcentered\ \CourseCode%
        \hfill \CourseProgram%
    \end{beamercolorbox}%
}

% Footer: minggu dan nomor halaman
\setbeamertemplate{footline}{%
    \begin{beamercolorbox}[wd=\paperwidth,ht=2.4ex,dp=1.2ex,leftskip=1em,rightskip=1em]{palette secondary}%
        \usebeamerfont{footline}\WeekLabel\ \textperiodcentered\ \WeekTopic%
        \hfill \insertframenumber\,/\,\inserttotalframenumber%
    \end{beamercolorbox}%
}

% Halaman judul: judul, subjudul, lalu blok identitas mata kuliah
\setbeamertemplate{title page}{%
    \begin{centering}
        \vskip1.0em
        {\usebeamerfont{title}\usebeamercolor[fg]{title}\inserttitle\par}
        \vskip0.3em
        {\usebeamerfont{subtitle}\usebeamercolor[fg]{subtitle}\insertsubtitle\par}
        \vskip1.2em
        \textcolor{ubbblue}{\rule{0.72\textwidth}{0.8pt}}
        \vskip1.2em
        {\small
            \textbf{\CourseName}\ (\CourseCode)\par
            \vskip0.4em
            \CourseProgram\par
            \CourseFaculty\par
            \CourseInstitution\par
            \vskip0.6em
            \CourseSemester\ \textperiodcentered\ \WeekLabel\par
            \vskip0.8em
            Dosen Pengampu: \CourseLecturer\par
        }
        \vskip1.0em
        \textcolor{ubbblue}{\rule{0.72\textwidth}{0.4pt}}
    \end{centering}
}

% =====================================================================
%  PERINTAH BANTU
% =====================================================================
\newcommand{\R}{\mathbb{R}}
\newcommand{\N}{\mathbb{N}}
\newcommand{\Z}{\mathbb{Z}}
\newcommand{\code}[1]{\texttt{#1}}

% Rujukan gambar dari buku: \bookfig[lebar]{NN}{namaberkas.png}
% Jalur relatif terhadap folder ppt/ (dikompilasi dari sana).
\newcommand{\bookfig}[3][0.80\linewidth]{%
    \includegraphics[width=#1]{../buku_ajar/chapters/#2/figures/#3}%
}

% Keterangan sumber gambar, sepadan dengan \sumber di buku
\newcommand{\figsumber}[1]{%
    \par\vspace{1pt}{\tiny\itshape Sumber: #1\par}%
}

% Penanda blok materi yang sering dipakai di slide
\newcommand{\slidejudul}[1]{{\large\bfseries\textcolor{ubbblue}{#1}}\par\vspace{2pt}}
\newcommand{\kunci}[1]{\textcolor{ubbred}{\textbf{#1}}}

%      \title{...} \subtitle{...} \setWeekInfo{Minggu n}{Topik}
%
%  Berkas minggu hanya berisi \setWeekInfo dan lingkungan frame.
%  Kompilasi dijalankan DARI DALAM folder ppt/ (lihat compile-ppt.bat),
%  karena gambar buku dirujuk dengan jalur relatif ../buku_ajar/.
% =====================================================================

% ---------------------------------------------------------------- bahasa
\usepackage[utf8]{inputenc}
\usepackage[T1]{fontenc}
\usepackage[indonesian]{babel}
\usepackage{lmodern}
\usepackage{microtype}

% ------------------------------------------------------------- matematika
\usepackage{amsmath,amssymb,mathtools}

% ----------------------------------------------------------------- tabel
\usepackage{booktabs}
\usepackage{array}
\usepackage{tabularx}

% --------------------------------------------------------------- grafik
\usepackage{tikz}
\usetikzlibrary{arrows.meta,positioning,shapes.geometric,calc,fit,backgrounds}

% ---------------------------------------------------------- kode program
\usepackage{listings}
\lstdefinestyle{slidecode}{
    basicstyle=\ttfamily\scriptsize,
    keywordstyle=\color{ubbblue}\bfseries,
    commentstyle=\color{ubbgreen}\itshape,
    stringstyle=\color{ubbred},
    numbers=left,
    numberstyle=\tiny\color{gray},
    numbersep=6pt,
    frame=single,
    rulecolor=\color{ubbblue!40},
    backgroundcolor=\color{ubbblue!3},
    breaklines=true,
    showstringspaces=false,
    tabsize=4,
    columns=fullflexible,
    keepspaces=true,
}
\lstset{style=slidecode}

% =====================================================================
%  WARNA INSTITUSI
% =====================================================================
\definecolor{ubbblue}{RGB}{12, 62, 120}      % warna utama
\definecolor{ubbteal}{RGB}{0, 118, 110}      % aksen kedua
\definecolor{ubbgreen}{RGB}{22, 122, 74}     % untuk blok contoh
\definecolor{ubbred}{RGB}{168, 42, 42}       % untuk blok peringatan
\definecolor{ubbgray}{RGB}{70, 76, 84}       % teks sekunder

% =====================================================================
%  IDENTITAS MATA KULIAH
% =====================================================================
\newcommand{\CourseName}{Pengantar Kriptografi}
\newcommand{\CourseCode}{MPFA014}
\newcommand{\CourseSemester}{Semester Ganjil}
\newcommand{\CourseProgram}{Program Studi Matematika}
\newcommand{\CourseFaculty}{Fakultas Matematika dan Ilmu Pengetahuan Alam}
\newcommand{\CourseInstitution}{Universitas Bale Bandung}
\newcommand{\CourseLecturer}{Cecep Suwanda, S.Si., M.Kom.}

% Label minggu, disetel tiap berkas lewat \setWeekInfo
\newcommand{\WeekLabel}{}
\newcommand{\WeekTopic}{}
\newcommand{\setWeekInfo}[2]{%
    \renewcommand{\WeekLabel}{#1}%
    \renewcommand{\WeekTopic}{#2}%
}

% =====================================================================
%  TEMA: header, footer, warna blok
% =====================================================================
\usetheme{default}
\setbeamertemplate{navigation symbols}{}
\setbeamertemplate{blocks}[rounded][shadow=false]

\setbeamercolor{palette primary}{bg=ubbblue, fg=white}
\setbeamercolor{palette secondary}{bg=ubbgray, fg=white}
\setbeamercolor{palette tertiary}{bg=ubbteal, fg=white}

\setbeamercolor{structure}{fg=ubbblue}
\setbeamercolor{normal text}{fg=black!92}
\setbeamercolor{alerted text}{fg=ubbred}

\setbeamercolor{block title}{bg=ubbblue, fg=white}
\setbeamercolor{block body}{bg=ubbblue!6, fg=black!92}
\setbeamercolor{block title alerted}{bg=ubbred, fg=white}
\setbeamercolor{block body alerted}{bg=ubbred!6, fg=black!92}
\setbeamercolor{block title example}{bg=ubbgreen, fg=white}
\setbeamercolor{block body example}{bg=ubbgreen!6, fg=black!92}

\setbeamercolor{itemize item}{fg=ubbblue}
\setbeamercolor{itemize subitem}{fg=ubbteal}
\setbeamercolor{enumerate item}{fg=ubbblue}

\setbeamerfont{headline}{size=\scriptsize}
\setbeamerfont{footline}{size=\scriptsize}
\setbeamerfont{block title}{size=\normalsize}

% Badan blok rata kiri (bukan rata kanan-kiri): teks slide lebih mudah dibaca
% dan baris terakhir yang pendek tidak lagi memicu "Underfull \hbox".
\addtobeamertemplate{block body begin}{\raggedright}{}
\addtobeamertemplate{block alerted body begin}{\raggedright}{}
\addtobeamertemplate{block example body begin}{\raggedright}{}

% Header: nama mata kuliah dan kode MK
\setbeamertemplate{headline}{%
    \begin{beamercolorbox}[wd=\paperwidth,ht=2.6ex,dp=1.2ex,leftskip=1em,rightskip=1em]{palette primary}%
        \usebeamerfont{headline}\CourseName\ \textperiodcentered\ \CourseCode%
        \hfill \CourseProgram%
    \end{beamercolorbox}%
}

% Footer: minggu dan nomor halaman
\setbeamertemplate{footline}{%
    \begin{beamercolorbox}[wd=\paperwidth,ht=2.4ex,dp=1.2ex,leftskip=1em,rightskip=1em]{palette secondary}%
        \usebeamerfont{footline}\WeekLabel\ \textperiodcentered\ \WeekTopic%
        \hfill \insertframenumber\,/\,\inserttotalframenumber%
    \end{beamercolorbox}%
}

% Halaman judul: judul, subjudul, lalu blok identitas mata kuliah
\setbeamertemplate{title page}{%
    \begin{centering}
        \vskip1.0em
        {\usebeamerfont{title}\usebeamercolor[fg]{title}\inserttitle\par}
        \vskip0.3em
        {\usebeamerfont{subtitle}\usebeamercolor[fg]{subtitle}\insertsubtitle\par}
        \vskip1.2em
        \textcolor{ubbblue}{\rule{0.72\textwidth}{0.8pt}}
        \vskip1.2em
        {\small
            \textbf{\CourseName}\ (\CourseCode)\par
            \vskip0.4em
            \CourseProgram\par
            \CourseFaculty\par
            \CourseInstitution\par
            \vskip0.6em
            \CourseSemester\ \textperiodcentered\ \WeekLabel\par
            \vskip0.8em
            Dosen Pengampu: \CourseLecturer\par
        }
        \vskip1.0em
        \textcolor{ubbblue}{\rule{0.72\textwidth}{0.4pt}}
    \end{centering}
}

% =====================================================================
%  PERINTAH BANTU
% =====================================================================
\newcommand{\R}{\mathbb{R}}
\newcommand{\N}{\mathbb{N}}
\newcommand{\Z}{\mathbb{Z}}
\newcommand{\code}[1]{\texttt{#1}}

% Rujukan gambar dari buku: \bookfig[lebar]{NN}{namaberkas.png}
% Jalur relatif terhadap folder ppt/ (dikompilasi dari sana).
\newcommand{\bookfig}[3][0.80\linewidth]{%
    \includegraphics[width=#1]{../buku_ajar/chapters/#2/figures/#3}%
}

% Keterangan sumber gambar, sepadan dengan \sumber di buku
\newcommand{\figsumber}[1]{%
    \par\vspace{1pt}{\tiny\itshape Sumber: #1\par}%
}

% Penanda blok materi yang sering dipakai di slide
\newcommand{\slidejudul}[1]{{\large\bfseries\textcolor{ubbblue}{#1}}\par\vspace{2pt}}
\newcommand{\kunci}[1]{\textcolor{ubbred}{\textbf{#1}}}


\title{Cipher Blok: Dari Satu Blok ke DES}
\subtitle{Mode operasi, jaringan Feistel, DES, dan Triple DES}
\setWeekInfo{Minggu 9}{Block cipher, DES, dan Triple DES}

\begin{document}

% ----------------------------------------------------------- halaman judul
\begin{frame}[plain]
    \titlepage
\end{frame}

% ------------------------------------------------------------------ peta
\begin{frame}{Peta Materi Pertemuan Ini}
    \footnotesize
    \begin{columns}[T,onlytextwidth]
        \begin{column}{0.52\textwidth}
            \slidejudul{Delapan bagian}
            \begin{enumerate}
                \item Cipher blok dan padding
                \item Lima mode operasi
                \item Confusion dan diffusion
                \item Iterated cipher dan Feistel
                \item DES: alur dan fungsi $f$
                \item Keamanan DES
                \item Double DES dan Triple DES
                \item Contoh hitung satu putaran
            \end{enumerate}
        \end{column}
        \begin{column}{0.44\textwidth}
            \begin{block}{Referensi}
                \footnotesize
                Bab 6 dan Bab 7 buku ajar; Munir, \textit{IF4020
                Kriptografi}, ITB, 2026 (modul 08 dan 09).
            \end{block}
            \begin{exampleblock}{Dua poros pertemuan}
                \footnotesize
                Bagaimana satu blok diacak, dan bagaimana blok-blok itu
                dirangkai menjadi pesan.
            \end{exampleblock}
        \end{column}
    \end{columns}
\end{frame}

% ---------------------------------------------------- capaian pembelajaran
\begin{frame}{Capaian Pembelajaran}
    \begin{block}{Sub-CPMK 1.3}
        \footnotesize
        Mengenal block cipher dan mode operasinya (ECB, CBC, CFB, OFB,
        CTR), jaringan Feistel, DES, Triple DES, AES, serta beberapa block
        cipher lain (GOST, RC5, RC6, Blowfish, Serpent, Twofish, dan MARS).
    \end{block}
    \vspace{0.1cm}
    {\scriptsize\itshape Pertemuan ini berhenti pada Triple DES; AES dan
    cipher blok lain menjadi bahan Minggu 10.}
    \vspace{0.15cm}
    \begin{exampleblock}{Setelah pertemuan ini, mahasiswa mampu}
        \footnotesize
        \begin{itemize}
            \item membedakan cipher blok dari cipher alir;
            \item menjelaskan peran lima mode operasi pada rangkaian blok;
            \item membedakan confusion dan diffusion;
            \item menelusuri satu putaran Feistel dan fungsi $f$ DES;
            \item menjelaskan mengapa Double DES lemah dan bagaimana
                  Triple DES menutupnya;
            \item menghitung putaran pertama DES pada vektor uji buku.
        \end{itemize}
    \end{exampleblock}
\end{frame}

% =====================================================================
\section{Cipher Blok: Satu Blok, Satu Kunci}
% =====================================================================

\begin{frame}{Cipher Blok Bekerja per Blok}
    \begin{columns}[T,onlytextwidth]
        \begin{column}{0.54\textwidth}
            \footnotesize
            \begin{itemize}
                \item Plainteks dipotong menjadi blok $n$ bit yang
                      panjangnya sama.
                \item Setiap blok dienkripsi sebagai satu kesatuan dengan
                      bit-bit kunci.
                \item Panjang cipherteks sama dengan panjang plainteks.
                \item Ukuran blok yang lazim: 64, 128, dan 256 bit.
                \item Panjang kunci eksternal tidak harus sama dengan
                      ukuran blok.
            \end{itemize}
        \end{column}
        \begin{column}{0.42\textwidth}
            \begin{block}{Dua fungsi}
                \[
                    C = E_K(P), \qquad P = D_K(C)
                \]
                \footnotesize
                Keduanya bekerja pada himpunan blok $n$ bit yang sama.
            \end{block}
            \begin{alertblock}{Beda dengan cipher alir}
                \footnotesize
                Cipher alir memproses data sedikit demi sedikit; cipher
                blok menunggu satu blok penuh.
            \end{alertblock}
        \end{column}
    \end{columns}
\end{frame}

\begin{frame}{Padding: Menutup Blok Terakhir}
    \begin{columns}[T,onlytextwidth]
        \begin{column}{0.50\textwidth}
            \footnotesize
            \begin{itemize}
                \item Pesan jarang habis dibagi ukuran blok.
                \item Blok terakhir yang pendek ditambahi bit
                      \textit{padding}.
                \item Isi padding dapat berupa bit 0 semua, bit 1 semua,
                      atau 0 dan 1 berselang-seling.
                \item Penerima harus tahu aturan padding agar dapat
                      membuangnya kembali.
            \end{itemize}
        \end{column}
        \begin{column}{0.46\textwidth}
            \begin{exampleblock}{Contoh blok 8 bit}
                \footnotesize
                Pesan tiga blok:
                \[
                    \code{10001101}\quad \code{00101001}\quad
                    \code{10110}
                \]
                Blok terakhir hanya 5 bit, ditambah tiga bit nol:
                \[
                    \code{10001101}\quad \code{00101001}\quad
                    \code{10110000}
                \]
            \end{exampleblock}
        \end{column}
    \end{columns}
    \vspace{0.15cm}
    \begin{block}{Ukuran blok dan ukuran kunci}
        \footnotesize
        Ukuran blok menentukan jangkauan diffusion; ukuran kunci
        menentukan efek confusion dan ketahanan terhadap pencarian
        menyeluruh.
    \end{block}
\end{frame}

\begin{frame}{Cipher Blok Mainan}
    \footnotesize
    Dua operasi sederhana sudah membentuk cipher blok yang dapat dibalik:
    XOR dengan kunci, lalu geser siklik satu bit ke kiri.
    \[
        C = E_K(P) = \mathrm{rotl}_1(P \oplus K), \qquad
        P = D_K(C) = \mathrm{rotr}_1(C) \oplus K
    \]
    \begin{columns}[T,onlytextwidth]
        \begin{column}{0.48\textwidth}
            \begin{exampleblock}{Enkripsi blok pertama}
                \footnotesize
                $P_1=\code{11011101}$, $K=\code{01010110}$:
                \begin{enumerate}
                    \item $P_1 \oplus K = \code{10001011}$;
                    \item geser kiri satu bit: $C_1=\code{00010111}$.
                \end{enumerate}
            \end{exampleblock}
        \end{column}
        \begin{column}{0.48\textwidth}
            \begin{exampleblock}{Dekripsi}
                \footnotesize
                \begin{enumerate}
                    \item Geser kanan $\code{00010111}$ menjadi
                          $\code{10001011}$;
                    \item XOR dengan kunci:
                          $\code{10001011}\oplus\code{01010110}
                          =\code{11011101}$.
                \end{enumerate}
            \end{exampleblock}
        \end{column}
    \end{columns}
    \vspace{0.1cm}
    \begin{alertblock}{Terlalu lemah}
        \footnotesize
        Cipher mainan ini mudah dipecahkan. Cipher blok modern mengulang
        substitusi, permutasi, pergeseran, kompresi, dan ekspansi berkali-kali.
    \end{alertblock}
\end{frame}

% =====================================================================
\section{Lima Mode Operasi}
% =====================================================================

\begin{frame}{Mengapa Perlu Mode Operasi}
    \footnotesize
    Fungsi $E_K$ hanya mengolah satu blok. Pesan yang dikirim biasanya
    terdiri atas banyak blok.
    \begin{block}{Definisi mode operasi}
        Aturan yang menentukan bagaimana blok-blok pesan dirangkai
        sebelum dienkripsi atau sesudah didekripsi.
    \end{block}
    \begin{columns}[T,onlytextwidth]
        \begin{column}{0.46\textwidth}
            \begin{enumerate}
                \item ECB --- \textit{Electronic Code Book}
                \item CBC --- \textit{Cipher Block Chaining}
                \item CFB --- \textit{Cipher Feedback}
                \item OFB --- \textit{Output Feedback}
                \item CTR --- \textit{Counter Mode}
            \end{enumerate}
        \end{column}
        \begin{column}{0.50\textwidth}
            \begin{exampleblock}{Dua pertanyaan pembeda}
                \footnotesize
                Apakah blok plainteks yang sama selalu menjadi blok
                cipherteks yang sama? Apakah galat satu bit menjalar ke
                blok lain?
            \end{exampleblock}
        \end{column}
    \end{columns}
\end{frame}

\begin{frame}{ECB: Tiap Blok Berdiri Sendiri}
    \footnotesize
    \[
        C_i = E_K(P_i), \qquad P_i = D_K(C_i)
    \]
    Tidak ada ketergantungan antarblok. Blok dapat diproses dalam urutan
    apa pun.
    \begin{exampleblock}{Contoh buku: plainteks \code{A23A9}, kunci \code{B}}
        \footnotesize
        Plainteks 20 bit dipotong menjadi lima blok 4 bit. Cipher blok
        mainan $\mathrm{rotl}_1(P\oplus K)$ diterapkan pada tiap blok.
        \begin{center}
            \begin{tabular}{@{}lccccc@{}}
                \toprule
                Blok & 1 & 2 & 3 & 4 & 5 \\
                \midrule
                $P_i$ & \code{1010} & \code{0010} & \code{0011}
                      & \code{1010} & \code{1001} \\
                $C_i$ & \code{0010} & \code{0011} & \code{0001}
                      & \code{0010} & \code{0100} \\
                \bottomrule
            \end{tabular}
        \end{center}
        Cipherteksnya \code{0010 0011 0001 0010 0100}, yaitu
        \code{23124} dalam heksadesimal.
    \end{exampleblock}
\end{frame}

\begin{frame}{ECB: Buku Kode dan Kebocoran Pola}
    \begin{columns}[T,onlytextwidth]
        \begin{column}{0.50\textwidth}
            \footnotesize
            \begin{itemize}
                \item Blok plainteks yang sama selalu menghasilkan blok
                      cipherteks yang sama.
                \item Pola berulang, misalnya spasi panjang, tetap tampak
                      dan dapat dianalisis frekuensinya.
                \item Secara teori mode ini adalah buku kode: $2^n$ buku
                      untuk kunci $n$ bit, masing-masing $2^n$ entri.
                \item Untuk blok 64 bit, satu buku berisi $2^{64}$ entri;
                      terlalu besar untuk disimpan.
            \end{itemize}
        \end{column}
        \begin{column}{0.46\textwidth}
            \begin{alertblock}{Blok dapat dibuang}
                \footnotesize
                Pesan ``Uang ditransfer lima satu juta rupiah'' dengan
                blok 2 karakter. Bila $C_8$ dan $C_9$ dibuang, penerima
                tetap memperoleh pesan bermakna: uang yang dikirim hanya
                satu juta rupiah.
            \end{alertblock}
            \begin{exampleblock}{Sisi baiknya}
                \footnotesize
                Akses acak dan galat satu bit hanya mengenai bloknya
                sendiri.
            \end{exampleblock}
        \end{column}
    \end{columns}
\end{frame}

\begin{frame}{CBC: Merantai Blok dengan IV}
    \begin{columns}[T,onlytextwidth]
        \begin{column}{0.54\textwidth}
            \footnotesize
            \[
                C_0 = \mathrm{IV}, \qquad
                C_i = E_K(P_i \oplus C_{i-1})
            \]
            \[
                P_i = D_K(C_i) \oplus C_{i-1}
            \]
            \begin{itemize}
                \item IV adalah blok semu pertama, diberikan pengguna
                      atau dibangkitkan acak.
                \item Setiap blok terikat pada seluruh blok sebelumnya.
                \item Blok plainteks yang sama pada posisi berbeda
                      menghasilkan cipherteks berbeda.
            \end{itemize}
        \end{column}
        \begin{column}{0.42\textwidth}
            \begin{center}
                \begin{tikzpicture}[>=Stealth,font=\scriptsize,
                                    scale=0.86,transform shape]
                    \node (p1) at (0,0) {$P_1$};
                    \node[circle,draw,inner sep=1.2pt] (x1) at (1.1,0)
                          {$\oplus$};
                    \node[draw,rounded corners,inner sep=2pt] (e1) at
                          (2.5,0) {$E_K$};
                    \node (c1) at (3.8,0) {$C_1$};
                    \node (iv) at (1.1,0.95) {IV};
                    \draw[->] (iv) -- (x1);
                    \draw[->] (p1) -- (x1);
                    \draw[->] (x1) -- (e1);
                    \draw[->] (e1) -- (c1);
                    \node (p2) at (0,-1.5) {$P_2$};
                    \node[circle,draw,inner sep=1.2pt] (x2) at (1.1,-1.5)
                          {$\oplus$};
                    \node[draw,rounded corners,inner sep=2pt] (e2) at
                          (2.5,-1.5) {$E_K$};
                    \node (c2) at (3.8,-1.5) {$C_2$};
                    \draw[->] (p2) -- (x2);
                    \draw[->] (c1) -- ++(0,-0.35) -| (x2);
                    \draw[->] (x2) -- (e2);
                    \draw[->] (e2) -- (c2);
                \end{tikzpicture}
            \end{center}
            \figsumber{Alur perantaian mode CBC.}
        \end{column}
    \end{columns}
\end{frame}

\begin{frame}{CBC: Pola Hilang, Galat Menjalar}
    \begin{columns}[T,onlytextwidth]
        \begin{column}{0.50\textwidth}
            \begin{exampleblock}{Contoh buku}
                \footnotesize
                Dengan plainteks \code{A23A9}, $K=\code{1011}$, dan
                $C_0=\code{0000}$, cipherteks CBC adalah
                \code{0010 0111 1111 1101 1111} atau \code{27FDF}.
                Blok \code{1010} di posisi 1 dan 4 menjadi \code{0010}
                dan \code{1101}, bukan lagi kembar.
            \end{exampleblock}
        \end{column}
        \begin{column}{0.46\textwidth}
            \footnotesize
            \begin{itemize}
                \item Pengulangan plainteks tidak lagi terlihat pada
                      cipherteks.
                \item Saat enkripsi, galat satu bit merambat ke blok itu
                      dan semua blok sesudahnya.
                \item Saat dekripsi, galat satu bit cipherteks hanya
                      mengenai blok itu dan satu bit pada blok berikutnya.
            \end{itemize}
        \end{column}
    \end{columns}
    \vspace{0.15cm}
    \begin{block}{Harga yang dibayar}
        \footnotesize
        Keamanan pola diperoleh dengan kehilangan sebagian sifat
        independen: enkripsi harus berjalan berurutan.
    \end{block}
\end{frame}

\begin{frame}{CFB: Cipher Blok Meniru Cipher Alir}
    \footnotesize
    Data dienkripsi per unit $s$ bit, dengan $s\le b$ dan $b$ ukuran blok.
    \[
        C_j = P_j \oplus \mathrm{MSB}_s(E_K(X_j)), \qquad
        X_{j+1} = \mathrm{LSB}_{b-s}(X_j) \parallel C_j
    \]
    \begin{columns}[T,onlytextwidth]
        \begin{column}{0.50\textwidth}
            \begin{itemize}
                \item $X_1=\mathrm{IV}$; antrian (\textit{shift register})
                      panjangnya $b$ bit.
                \item Umpan balik diambil dari cipherteks.
                \item Bila $s=b$: $C_i=P_i\oplus E_K(C_{i-1})$.
                \item Dekripsi hanya memerlukan $E_K$, bukan $D_K$.
            \end{itemize}
        \end{column}
        \begin{column}{0.46\textwidth}
            \begin{alertblock}{Galat}
                \footnotesize
                Pada enkripsi, galat merambat ke blok sekarang dan
                blok-blok berikutnya. Pada dekripsi, galat terbatas pada
                blok sekarang dan satu blok sesudahnya.
            \end{alertblock}
            \begin{exampleblock}{Kegunaan}
                \footnotesize
                Cocok untuk data yang belum mencapai satu blok penuh,
                misalnya satu karakter 8 bit (CFB 8-bit).
            \end{exampleblock}
        \end{column}
    \end{columns}
\end{frame}

\begin{frame}{OFB: Umpan Balik dari Keluaran Enkripsi}
    \footnotesize
    \[
        C_j = P_j \oplus \mathrm{MSB}_s(E_K(X_j)), \qquad
        X_{j+1} = \mathrm{LSB}_{b-s}(X_j) \parallel
        \mathrm{MSB}_s(E_K(X_j))
    \]
    \begin{columns}[T,onlytextwidth]
        \begin{column}{0.50\textwidth}
            \begin{itemize}
                \item Bedanya dengan CFB: yang masuk antrian adalah
                      keluaran $E_K$, bukan cipherteks.
                \item Umpan balik tidak melewati cipherteks.
                \item Galat satu bit pada plainteks maupun cipherteks
                      hanya mengenai bloknya sendiri.
            \end{itemize}
        \end{column}
        \begin{column}{0.46\textwidth}
            \begin{exampleblock}{Sasaran pemakaian}
                \footnotesize
                Transmisi analog yang didigitisasi, seperti suara dan
                video. Satu bit salah masih dapat ditoleransi, tetapi
                penjalaran galat tidak boleh terjadi.
            \end{exampleblock}
        \end{column}
    \end{columns}
\end{frame}

\begin{frame}{CTR: Counter sebagai Pembangkit Keystream}
    \footnotesize
    \[
        C_j = P_j \oplus E_K(T_j), \qquad T_j = T_{j-1}+1
    \]
    \begin{columns}[T,onlytextwidth]
        \begin{column}{0.50\textwidth}
            \begin{itemize}
                \item $T_j$ adalah blok counter sebesar ukuran blok.
                \item Nilai counter harus berbeda untuk setiap blok.
                \item Counter diinisialisasi sekali, lalu dinaikkan satu
                      untuk blok berikutnya.
                \item Blok terakhir yang tidak penuh memakai bit paling
                      kiri dari keystream.
            \end{itemize}
        \end{column}
        \begin{column}{0.46\textwidth}
            \begin{exampleblock}{Sifat utama}
                \footnotesize
                Hanya algoritma enkripsi yang diperlukan. Blok tidak
                saling bergantung, sehingga enkripsi dan dekripsi dapat
                diparalelkan dan diakses acak.
            \end{exampleblock}
            \begin{alertblock}{Syarat}
                \footnotesize
                Nilai counter harus tidak diketahui atau tidak dapat
                diprediksi lawan.
            \end{alertblock}
        \end{column}
    \end{columns}
\end{frame}

\begin{frame}{Perbandingan Lima Mode}
    \scriptsize
    \begin{center}
        \begin{tabularx}{\textwidth}{@{}>{\raggedright\arraybackslash}p{1.05cm}
            *{3}{>{\raggedright\arraybackslash}X}
            >{\raggedright\arraybackslash}p{3.3cm}@{}}
            \toprule
            \textbf{Mode} & \textbf{Ketergantungan} &
            \textbf{Galat} & \textbf{Akses} & \textbf{Ciri khas} \\
            \midrule
            ECB & Tidak ada & Hanya blok itu & Acak, paralel &
            Pola plainteks berulang ikut bocor \\
            CBC & Blok sebelumnya & Enkripsi menjalar & Enkripsi serial &
            IV menghapus pola berulang \\
            CFB & Cipherteks & Menjalar saat enkripsi & Serial &
            Bekerja per unit $s$ bit seperti cipher alir \\
            OFB & Keluaran $E_K$ & Tidak menjalar & Serial &
            Cocok untuk suara dan video terdigitalisasi \\
            CTR & Tidak ada & Tidak menjalar & Acak, paralel &
            Counter harus unik dan tak terprediksi \\
            \bottomrule
        \end{tabularx}
    \end{center}
    \vspace{0.1cm}
    {\footnotesize Pilihan mode menentukan dua hal sekaligus: pola apa
    yang masih terlihat pada cipherteks, dan sejauh mana satu galat
    merusak data sesudahnya.}
\end{frame}

% =====================================================================
\section{Confusion dan Diffusion}
% =====================================================================

\begin{frame}{Confusion: Menutup Hubungan Statistik}
    \begin{columns}[T,onlytextwidth]
        \begin{column}{0.52\textwidth}
            \footnotesize
            \begin{itemize}
                \item Diperkenalkan Shannon pada 1949 bersama diffusion.
                \item Confusion menyembunyikan hubungan statistik antara
                      plainteks, cipherteks, dan kunci.
                \item Kriptanalis kehilangan pola statistik yang biasa
                      dipakai untuk menebak kunci.
                \item Sarannya: substitusi nirlanjar, umumnya berupa tabel
                      pencarian.
            \end{itemize}
        \end{column}
        \begin{column}{0.44\textwidth}
            \begin{block}{Inti gagasan}
                \footnotesize
                Hubungan antara bit masukan dan bit keluaran dibuat
                berbelit, sehingga perubahan kecil pada kunci tidak
                tampak sebagai perubahan yang dapat diprediksi.
            \end{block}
        \end{column}
    \end{columns}
\end{frame}

\begin{frame}{S-box: Enam Bit Menjadi Empat Bit}
    \footnotesize
    Pada DES, kotak-S menerima 6 bit dan mengeluarkan 4 bit.
    \begin{columns}[T,onlytextwidth]
        \begin{column}{0.44\textwidth}
            \begin{itemize}
                \item Bit ke-1 dan ke-6 menentukan nomor baris.
                \item Bit ke-2 sampai ke-5 menentukan nomor kolom.
                \item Masukan $\code{110100}$: baris
                      $(10)_2=2$, kolom $(1010)_2=10$.
                \item Masukan $\code{101010}$: baris 2, kolom 5.
            \end{itemize}
        \end{column}
        \begin{column}{0.54\textwidth}
            \begin{exampleblock}{Baris 2 tabel $S_1$}
                \scriptsize
                \setlength{\tabcolsep}{1.2pt}
                \begin{center}
                    \begin{tabular}{@{}*{16}{c}@{}}
                        4 & 1 & 14 & 8 & 13 & 6 & 2 & 11 & 15 & 12
                        & 9 & 7 & 3 & 10 & 5 & 0 \\
                    \end{tabular}
                \end{center}
                Kolom ke-10 bernilai 9, jadi keluarannya
                $\code{1001}$.
            \end{exampleblock}
            \begin{alertblock}{Peran dalam cipher}
                \footnotesize
                Mengubah satu bit masukan mengubah paling sedikit dua bit
                keluaran. Sifat ini menopang confusion.
            \end{alertblock}
        \end{column}
    \end{columns}
\end{frame}

\begin{frame}{Diffusion: Satu Bit Mengubah Banyak Bit}
    \footnotesize
    Diffusion menyebarkan pengaruh satu bit plainteks atau kunci ke
    sebanyak mungkin bit cipherteks.
    \begin{exampleblock}{Contoh blok 16 bit}
        \footnotesize
        \[
            \code{1111111111111111} \longrightarrow
            \code{0110110000010100}
        \]
        Ubah bit pertama menjadi 0:
        \[
            \code{0111111111111111} \longrightarrow
            \code{1011001011101111}
        \]
        Perubahan satu bit masukan mengubah banyak bit keluaran.
    \end{exampleblock}
    \begin{columns}[T,onlytextwidth]
        \begin{column}{0.48\textwidth}
            \begin{itemize}
                \item Dicapai dengan permutasi atau transposisi yang
                      diulang.
                \item Hasilnya sukar diprediksi dari masukan.
            \end{itemize}
        \end{column}
        \begin{column}{0.48\textwidth}
            \begin{block}{Akibat bagi penyerang}
                \footnotesize
                Menebak kunci dari sebagian cipherteks tidak lagi
                memberi petunjuk tentang bit plainteks yang lain.
            \end{block}
        \end{column}
    \end{columns}
\end{frame}

\begin{frame}{P-box dan Putaran yang Diulang}
    \begin{columns}[T,onlytextwidth]
        \begin{column}{0.54\textwidth}
            \footnotesize
            \begin{itemize}
                \item P-box adalah matriks permutasi: entri ke-$i$
                      menyatakan bit masukan yang pindah ke posisi $i$.
                \item Pada DES, bit ke-58 pindah ke posisi pertama dan
                      bit ke-50 ke posisi kedua.
                \item Permutasi dapat pula mengompresi (mengurangi bit)
                      atau mengekspansi (menambah bit).
                \item Cipher lain memakai pergeseran siklik, bukan P-box.
            \end{itemize}
        \end{column}
        \begin{column}{0.42\textwidth}
            \begin{block}{Satu putaran}
                \footnotesize
                Substitusi menghasilkan confusion, permutasi
                menghasilkan diffusion. Keduanya diulang dengan kunci
                putaran yang berbeda.
            \end{block}
        \end{column}
    \end{columns}
\end{frame}

% =====================================================================
\section{Iterated Cipher dan Jaringan Feistel}
% =====================================================================

\begin{frame}{Iterated Cipher: Satu Blok, $r$ Putaran}
    \footnotesize
    \[
        C_i = f(C_{i-1}, K_i), \qquad i=1,2,\dots,r
    \]
    Plainteks adalah $C_0$, cipherteks adalah $C_r$. Kunci sering
    dicampurkan sebelum putaran pertama dan sesudah putaran terakhir:
    \[
        C_0 = P\oplus K_0, \qquad
        C = C_{r-1}\oplus K_r
    \]
    \begin{columns}[T,onlytextwidth]
        \begin{column}{0.50\textwidth}
            \begin{itemize}
                \item $f$ disebut fungsi putaran; di dalamnya ada
                      substitusi dan permutasi.
                \item Setiap putaran memakai subkunci berbeda.
                \item DES memakai $r=16$.
            \end{itemize}
        \end{column}
        \begin{column}{0.46\textwidth}
            \begin{alertblock}{Trade-off}
                \footnotesize
                Makin banyak putaran, makin luas diffusion, tetapi
                komputasi makin mahal. Jumlah putaran harus
                menyeimbangkan keamanan dan efisiensi.
            \end{alertblock}
        \end{column}
    \end{columns}
\end{frame}

\begin{frame}{Jaringan Feistel: Dua Separuh yang Bertukar Peran}
    \begin{columns}[T,onlytextwidth]
        \begin{column}{0.52\textwidth}
            \footnotesize
            Blok dibagi menjadi $L$ dan $R$. Pada setiap putaran:
            \[
                L_i = R_{i-1}, \qquad
                R_i = L_{i-1}\oplus f(R_{i-1},K_i)
            \]
            \begin{itemize}
                \item Hanya satu separuh yang diolah fungsi $f$.
                \item Hasilnya di-XOR-kan ke separuh lain.
                \item Kedua separuh bertukar peran pada putaran
                      berikutnya.
                \item Struktur ini dipakai DES.
            \end{itemize}
        \end{column}
        \begin{column}{0.44\textwidth}
            \begin{center}
                \bookfig[0.98\linewidth]{06}{jaringan-feistel.png}
            \end{center}
            \figsumber{Munir, \textit{IF4020 Kriptografi}, ITB.}
        \end{column}
    \end{columns}
\end{frame}

\begin{frame}{Mengapa Feistel Mudah Dibalik}
    \footnotesize
    Dari persamaan enkripsi, dekripsi tidak memerlukan fungsi $f^{-1}$:
    \[
        R_{i-1} = L_i, \qquad
        L_{i-1} = R_i \oplus f(L_i,K_i)
    \]
    Pemeriksaannya langsung:
    \[
        L_{i-1}\oplus f(R_{i-1},K_i)\oplus f(R_{i-1},K_i)
        = L_{i-1}
    \]
    \begin{columns}[T,onlytextwidth]
        \begin{column}{0.52\textwidth}
            \begin{itemize}
                \item Dekripsi menjalankan jaringan yang sama dari bawah
                      ke atas.
                \item Urutan subkunci dibalik: $K_{16},\dots,K_1$.
                \item Perancang tidak perlu menulis algoritma dekripsi
                      terpisah.
            \end{itemize}
        \end{column}
        \begin{column}{0.44\textwidth}
            \begin{exampleblock}{Konsekuensi rancangan}
                \footnotesize
                Karena tidak perlu dibalik, $f$ boleh dibuat serumit
                mungkin tanpa mengorbankan kemampuan dekripsi.
            \end{exampleblock}
        \end{column}
    \end{columns}
\end{frame}

\begin{frame}{Key Scheduling: Dari Kunci Eksternal ke Subkunci}
    \footnotesize
    \begin{columns}[T,onlytextwidth]
        \begin{column}{0.50\textwidth}
            \begin{itemize}
                \item Kunci eksternal diturunkan menjadi kunci putaran
                      melalui \textit{key expansion} atau
                      \textit{key scheduling}.
                \item Panjang subkunci tidak harus sama dengan panjang
                      kunci eksternal.
                \item DES: kunci 64 bit menjadi 16 subkunci 48 bit.
                \item Subkunci berbeda tiap putaran mencegah transformasi
                      yang sama terulang.
            \end{itemize}
        \end{column}
        \begin{column}{0.46\textwidth}
            \begin{center}
                \bookfig[0.98\linewidth]{06}{pembangkitan-subkunci.png}
            \end{center}
            \figsumber{Munir, \textit{IF4020 Kriptografi}, ITB.}
        \end{column}
    \end{columns}
\end{frame}

% =====================================================================
\section{Data Encryption Standard}
% =====================================================================

\begin{frame}{DES: Standar, Algoritma, dan Sejarah}
    \begin{columns}[T,onlytextwidth]
        \begin{column}{0.52\textwidth}
            \footnotesize
            \begin{itemize}
                \item DES adalah nama standar; algoritmanya disebut
                      \textit{Data Encryption Algorithm} (DEA).
                \item Termasuk cipher blok kunci-simetri.
                \item 1972: NBS meminta standar enkripsi blok.
                \item 1974: Horst Feistel di IBM merancang
                      \textit{Lucifer} dengan kunci dan blok 128 bit.
                \item Lucifer berkembang menjadi DES: kunci 56 bit, blok
                      64 bit, 16 putaran.
                \item 1976: DES disetujui setelah dinilai NSA.
            \end{itemize}
        \end{column}
        \begin{column}{0.44\textwidth}
            \begin{block}{Yang berubah dari Lucifer}
                \footnotesize
                Panjang kunci dipangkas, ukuran blok diperkecil, dan
                jumlah putaran ditetapkan 16. Perubahan itulah yang
                menjadi inti perdebatan keamanan DES.
            \end{block}
        \end{column}
    \end{columns}
\end{frame}

\begin{frame}{Parameter DES}
    \begin{columns}[T,onlytextwidth]
        \begin{column}{0.56\textwidth}
            \small
            \begin{tabular}{@{}ll@{}}
                \toprule
                \textbf{Parameter} & \textbf{Nilai} \\
                \midrule
                Ukuran blok & 64 bit \\
                Kunci eksternal & 64 bit \\
                Kunci efektif & 56 bit \\
                Bit paritas & 8, 16, 24, 32, 40, 48, 56, 64 \\
                Putaran & 16 \\
                Kunci internal & 16 buah, masing-masing 48 bit \\
                \bottomrule
            \end{tabular}
        \end{column}
        \begin{column}{0.40\textwidth}
            \begin{alertblock}{Delapan bit yang dibuang}
                \footnotesize
                Bit pada posisi kelipatan 8 hanya paritas. Karena itu
                ruang kunci DES berisi $2^{56}$ kemungkinan, bukan
                $2^{64}$.
            \end{alertblock}
        \end{column}
    \end{columns}
    \vspace{0.25cm}
    \begin{block}{Ruang kunci}
        \footnotesize
        Angka $2^{56}$ inilah yang menentukan keamanan DES dan menjadi
        sasaran serangan pencarian menyeluruh.
    \end{block}
\end{frame}

\begin{frame}{Alur Enkripsi DES}
    \begin{columns}[T,onlytextwidth]
        \begin{column}{0.50\textwidth}
            \footnotesize
            Tiga tahap:
            \begin{enumerate}
                \item Permutasi awal (IP) mengubah urutan 64 bit
                      plainteks.
                \item 16 putaran \textit{enciphering} memakai
                      $K_1,\dots,K_{16}$.
                \item Invers permutasi awal ($\mathrm{IP}^{-1}$)
                      membentuk cipherteks.
            \end{enumerate}
            \begin{itemize}
                \item Setelah putaran ke-16, blok disusun
                      $R_{16}\parallel L_{16}$; putaran terakhir tidak
                      menukar separuh.
                \item Dekripsi memakai algoritma sama dengan subkunci
                      terbalik.
            \end{itemize}
        \end{column}
        \begin{column}{0.46\textwidth}
            \begin{center}
                \bookfig[0.98\linewidth]{07}{skema-enkripsi-des.png}
            \end{center}
            \figsumber{Munir, \textit{IF4020 Kriptografi}, ITB.}
        \end{column}
    \end{columns}
\end{frame}

\begin{frame}{Fungsi $f$: Empat Langkah}
    \begin{columns}[T,onlytextwidth]
        \begin{column}{0.54\textwidth}
            \footnotesize
            Fungsi $f$ menerima 32 bit dan subkunci 48 bit, lalu
            mengembalikan 32 bit.
            \begin{enumerate}
                \item Ekspansi $E$: 32 bit menjadi 48 bit.
                \item XOR dengan $K_i$: $A=E(R_{i-1})\oplus K_i$.
                \item $A$ dipecah menjadi delapan kelompok 6 bit;
                      $S_1,\dots,S_8$ masing-masing menghasilkan 4 bit,
                      membentuk $B$ 32 bit.
                \item Permutasi $P$ mengacak $B$.
            \end{enumerate}
            \[
                f(R,K)=P\bigl(S_1(B_1)\parallel\dots\parallel
                S_8(B_8)\bigr)
            \]
            Luaran putaran: $(L_i,R_i)=(R_{i-1},L_{i-1}\oplus P(B))$.
        \end{column}
        \begin{column}{0.42\textwidth}
            \begin{center}
                \bookfig[0.98\linewidth]{07}{fungsi-f-des.png}
            \end{center}
            \figsumber{Munir, \textit{IF4020 Kriptografi}, ITB.}
        \end{column}
    \end{columns}
\end{frame}

\begin{frame}{Pembangkitan 16 Subkunci DES}
    \begin{columns}[T,onlytextwidth]
        \begin{column}{0.50\textwidth}
            \footnotesize
            \begin{enumerate}
                \item PC-1 membuang delapan bit paritas dan menyisakan
                      56 bit.
                \item Hasilnya dipecah menjadi $C_0$ dan $D_0$, masing-masing
                      28 bit.
                \item $C_i$ dan $D_i$ digeser siklik ke kiri menurut
                      jadwal $1,1,2,2,\dots,1$.
                \item PC-2 memilih 48 bit dari $C_i\parallel D_i$
                      menjadi $K_i$.
            \end{enumerate}
            \[
                \begin{aligned}
                    C_i&=\mathrm{LS}_{s_i}(C_{i-1}), \qquad
                    D_i=\mathrm{LS}_{s_i}(D_{i-1}),\\
                    K_i&=\mathrm{PC2}(C_i\parallel D_i)
                \end{aligned}
            \]
        \end{column}
        \begin{column}{0.46\textwidth}
            \begin{center}
                \bookfig[0.98\linewidth]{07}{pembangkitan-kunci-des.png}
            \end{center}
            \figsumber{Munir, \textit{IF4020 Kriptografi}, ITB.}
        \end{column}
    \end{columns}
\end{frame}

% =====================================================================
\section{Keamanan DES}
% =====================================================================

\begin{frame}{Ruang Kunci $2^{56}$ dan Pencarian Menyeluruh}
    \footnotesize
    \begin{columns}[T,onlytextwidth]
        \begin{column}{0.52\textwidth}
            \begin{itemize}
                \item Penyerang memerlukan pasangan plainteks-cipherteks
                      yang berkoresponden.
                \item Untuk setiap kandidat $k$, hitung
                      $\mathrm{DES}(k,p)$ dan bandingkan dengan $c$.
                \item Ruang kunci berisi
                      \[
                          2^{56}=72\,057\,594\,037\,927\,936
                      \]
                      kemungkinan.
                \item Kunci diharapkan ditemukan sebelum $2^{55}$
                      iterasi.
            \end{itemize}
        \end{column}
        \begin{column}{0.44\textwidth}
            \begin{exampleblock}{Laju satu juta kunci per detik}
                \footnotesize
                Pencarian rata-rata $2^{55}$ memerlukan sekitar 1.142
                tahun. Seluruh ruang kunci dua kali lebih lama.
            \end{exampleblock}
            \begin{alertblock}{Kenyataan perangkat keras}
                \footnotesize
                1997: kontes DESCHALL. 1998: mesin EFF senilai \$250.000
                dengan 1.856 chip menemukan kunci dalam sedikit lebih
                dari dua hari. 1999: kurang dari satu hari.
            \end{alertblock}
        \end{column}
    \end{columns}
\end{frame}

\begin{frame}{Efek Longsoran}
    \footnotesize
    Efek longsoran menunjukkan bahwa perubahan satu bit masukan menyebar
    ke sekitar separuh bit keluaran.
    \begin{exampleblock}{Vektor uji resmi}
        \footnotesize
        \[
            \mathrm{DES}(\code{0123456789ABCDEF},
            \code{133457799BBCDFF1})
            = \code{85E813540F0AB405}
        \]
        Balik bit pertama plainteks:
        \[
            \mathrm{DES}(\code{8123456789ABCDEF},
            \code{133457799BBCDFF1})
            = \code{AD9DFD6C35DE8DEC}
        \]
    \end{exampleblock}
    \begin{columns}[T,onlytextwidth]
        \begin{column}{0.48\textwidth}
            \begin{itemize}
                \item Kedua cipherteks berbeda pada 33 dari 64 bit.
                \item Perubahan itu mendekati separuh bit.
            \end{itemize}
        \end{column}
        \begin{column}{0.48\textwidth}
            \begin{block}{Makna praktis}
                \footnotesize
                Hubungan masukan-keluaran sukar ditebak, sehingga
                penyerang tidak dapat menilai kandidat kunci dari
                kemiripan cipherteks.
            \end{block}
        \end{column}
    \end{columns}
\end{frame}

\begin{frame}{Empat Kunci Lemah DES}
    \begin{columns}[T,onlytextwidth]
        \begin{column}{0.52\textwidth}
            \footnotesize
            \begin{itemize}
                \item Kunci lemah memenuhi $E_K(E_K(x))=x$ untuk semua
                      plainteks $x$.
                \item Artinya enkripsi dan dekripsi menjadi sama.
                \item Penyebabnya: seluruh 16 subkunci identik.
                \item Ada empat kunci lemah, semuanya panjang 56 bit.
            \end{itemize}
        \end{column}
        \begin{column}{0.44\textwidth}
            \begin{block}{Contoh \code{0101010101010101}}
                \footnotesize
                Setiap byte bernilai \code{00000001}. PC-1 membuang bit
                terakhir tiap byte, sehingga $C_0$ dan $D_0$ seluruhnya
                nol. Pergeseran atas deretan nol tetap nol, dan semua
                subkunci menjadi \code{000000000000}.
            \end{block}
        \end{column}
    \end{columns}
    \vspace{0.15cm}
    \begin{center}
        \small
        \code{0101010101010101} \quad
        \code{1F1F1F1F0E0E0E0E} \\
        \vspace{2pt}
        \code{E0E0E0E0F1F1F1F1} \quad
        \code{FEFEFEFEFEFEFEFE}
    \end{center}
    \begin{alertblock}{Praktik}
        \footnotesize
        Keempat kunci ini harus dihindari saat membangkitkan kunci DES.
    \end{alertblock}
\end{frame}

% =====================================================================
\section{Double DES dan Triple DES}
% =====================================================================

\begin{frame}{Double DES dan Serangan Meet-in-the-Middle}
    \footnotesize
    \[
        C=E_{K_2}\bigl(E_{K_1}(P)\bigr), \qquad
        P=D_{K_1}\bigl(D_{K_2}(C)\bigr)
    \]
    \begin{columns}[T,onlytextwidth]
        \begin{column}{0.52\textwidth}
            \begin{enumerate}
                \item Tentukan nilai antara
                      $X=E_{K_1}(P)=D_{K_2}(C)$.
                \item Enkripsi $P$ dengan semua $2^{56}$ nilai $K_1$;
                      simpan setiap $X$ dalam tabel.
                \item Dekripsi $C$ dengan semua $2^{56}$ nilai $K_2$;
                      bandingkan hasilnya dengan isi tabel.
                \item Kecocokan memberi pasangan calon $(K_1,K_2)$.
                \item Uji pasangan itu dengan pasangan
                      plainteks-cipherteks lain.
            \end{enumerate}
        \end{column}
        \begin{column}{0.44\textwidth}
            \begin{alertblock}{Mengapa gagal}
                \footnotesize
                Dua kunci 56 bit seolah memberi $2^{112}$ kemungkinan,
                tetapi pertemuan di nilai antara menurunkan biaya
                pencarian menjadi sekitar $2^{57}$.
            \end{alertblock}
            \begin{block}{Pelajaran}
                \footnotesize
                Menambah satu lapis enkripsi tidak otomatis menggandakan
                keamanan.
            \end{block}
        \end{column}
    \end{columns}
\end{frame}

\begin{frame}{Triple DES: EEE dan EDE}
    \begin{columns}[T,onlytextwidth]
        \begin{column}{0.54\textwidth}
            \footnotesize
            \textbf{Mode EEE} (enkripsi-enkripsi-enkripsi):
            \[
                C=E_{K_3}\bigl(E_{K_2}(E_{K_1}(P))\bigr)
            \]
            \textbf{Mode EDE} (enkripsi-dekripsi-enkripsi):
            \[
                C=E_{K_3}\bigl(D_{K_2}(E_{K_1}(P))\bigr)
            \]
            \[
                P=D_{K_1}\bigl(E_{K_2}(D_{K_3}(C))\bigr)
            \]
        \end{column}
        \begin{column}{0.42\textwidth}
            \begin{block}{Mengapa EDE}
                \footnotesize
                Bila $K_1=K_2=K_3=K$, maka
                $E_K(D_K(E_K(P)))=E_K(P)$. Triple DES tetap kompatibel
                dengan DES tunggal.
            \end{block}
            \begin{exampleblock}{Panjang kunci}
                \footnotesize
                Dua kunci: $K_1,K_2,K_1$ memberi $56+56=112$ bit.
                Tiga kunci: $K_1,K_2,K_3$ memberi $56+56+56=168$ bit.
            \end{exampleblock}
        \end{column}
    \end{columns}
    \vspace{0.1cm}
    \begin{alertblock}{Tujuan}
        \footnotesize
        Tiga lapis enkripsi membuat serangan meet-in-the-middle kembali
        mahal, tanpa meninggalkan kompatibilitas dengan DES.
    \end{alertblock}
\end{frame}

% =====================================================================
\section{Contoh Terhitung: Putaran Pertama DES}
% =====================================================================

\begin{frame}{Contoh: Putaran Pertama DES (1/4)}
    \begin{columns}[T,onlytextwidth]
        \begin{column}{0.44\textwidth}
            \begin{block}{Vektor uji}
                \small
                \begin{tabular}{@{}ll@{}}
                    $P$ & \code{0123456789ABCDEF} \\
                    $K$ & \code{133457799BBCDFF1} \\
                    $C$ & \code{85E813540F0AB405} \\
                \end{tabular}
            \end{block}
            \begin{exampleblock}{Kunci pertama}
                \footnotesize
                Dari jadwal kunci diperoleh
                $K_1=\code{1B02EFFC7072}$.
            \end{exampleblock}
        \end{column}
        \begin{column}{0.52\textwidth}
            \textbf{Langkah 1 --- permutasi awal (IP).}
            \footnotesize
            Membaca tabel IP pada 64 bit plainteks menghasilkan dua
            separuh:
            \[
                L_0=\code{CC00CCFF}, \qquad
                R_0=\code{F0AAF0AA}
            \]
            Operasi: urutan bit plainteks diubah menurut tabel IP, lalu
            hasilnya dibelah dua.
            \[
                \mathrm{IP}(P)=\code{CC00CCFFF0AAF0AA}
            \]
        \end{column}
    \end{columns}
\end{frame}

\begin{frame}{Contoh: Putaran Pertama DES (2/4)}
    \footnotesize
    \textbf{Langkah 2 --- ekspansi dan XOR subkunci.}
    $R_0$ dalam biner:
    \[
        \code{1111}\ \code{0000}\ \code{1010}\ \code{1010}\
        \code{1111}\ \code{0000}\ \code{1010}\ \code{1010}
    \]
    Ekspansi $E$ mengubah 32 bit menjadi 48 bit:
    \[
        E(R_0)=\code{7A15557A1555}
    \]
    XOR dengan $K_1$:
    \[
        \begin{aligned}
            E(R_0)\oplus K_1
            &=\code{011000}\ \code{010001}\ \code{011110}\
              \code{111010}\\
            &\quad\ \code{100001}\ \code{100110}\ \code{010100}\
              \code{100111}\\
            &=\code{6117BA866527}
        \end{aligned}
    \]
    \begin{block}{Hasil langkah 2}
        Diperoleh blok $A=\code{6117BA866527}$ berukuran 48 bit, siap
        dipecah menjadi delapan kelompok enam bit.
    \end{block}
\end{frame}

\begin{frame}{Contoh: Putaran Pertama DES (3/4)}
    \footnotesize
    \textbf{Langkah 3 --- substitusi kotak-S.} Setiap kelompok enam bit
    masuk ke kotak-S yang berbeda.
    \begin{center}
        \small
        \begin{tabular}{@{}cccccc@{}}
            \toprule
            $j$ & Masukan & Baris & Kolom & $S_j$ & Keluaran \\
            \midrule
            1 & \code{011000} & 0 & 12 & 5  & \code{0101} \\
            2 & \code{010001} & 1 & 8  & 12 & \code{1100} \\
            3 & \code{011110} & 0 & 15 & 8  & \code{1000} \\
            4 & \code{111010} & 2 & 13 & 2  & \code{0010} \\
            5 & \code{100001} & 3 & 0  & 11 & \code{1011} \\
            6 & \code{100110} & 2 & 3  & 5  & \code{0101} \\
            7 & \code{010100} & 0 & 10 & 9  & \code{1001} \\
            8 & \code{100111} & 3 & 3  & 7  & \code{0111} \\
            \bottomrule
        \end{tabular}
    \end{center}
    \begin{block}{Hasil langkah 3}
        Delapan keluaran 4 bit digabung menjadi
        $S\text{-out}=\code{5C82B597}$ (32 bit). Setiap baris memakai
        kotak-S yang berbeda; hanya potongan tabel yang diperlukan
        contoh ini yang ditampilkan.
    \end{block}
\end{frame}

\begin{frame}{Contoh: Putaran Pertama DES (4/4)}
    \footnotesize
    \textbf{Langkah 4 --- permutasi $P$.}
    Membaca tabel $P$ pada $\code{5C82B597}$:
    \[
        P(S)=\code{234AA9BB}
    \]
    \textbf{Langkah 5 --- pembentukan separuh baru.}
    Subkunci $K_1$ tidak dipakai lagi; yang tersisa adalah XOR dengan
    separuh kiri:
    \[
        R_1=L_0\oplus P(S)
        =\code{CC00CCFF}\oplus\code{234AA9BB}
        =\code{EF4A6544}
    \]
    \[
        L_1=R_0=\code{F0AAF0AA}
    \]
    \begin{columns}[T,onlytextwidth]
        \begin{column}{0.48\textwidth}
            \begin{block}{Setelah 16 putaran}
                \footnotesize
                $L_{16}=\code{43423234}$ dan
                $R_{16}=\code{0A4CD995}$. Blok sebelum permutasi akhir:
                $\code{0A4CD99543423234}$.
            \end{block}
        \end{column}
        \begin{column}{0.48\textwidth}
            \begin{exampleblock}{Hasil akhir}
                \footnotesize
                $\mathrm{IP}^{-1}$ memberi
                $C=\code{85E813540F0AB405}$, sama dengan vektor uji.
                Dekripsi dengan $K_{16},\dots,K_1$ mengembalikan $P$.
            \end{exampleblock}
        \end{column}
    \end{columns}
\end{frame}

% ------------------------------------------------------------------ penutup
\begin{frame}{Penutup}
    \begin{block}{Yang harus dibawa dari pertemuan ini}
        \begin{enumerate}
            \item Cipher blok mengolah blok $n$ bit; mode operasi
                  menentukan perantaian, pola yang bocor, dan penjalaran
                  galat.
            \item ECB membocorkan pola; CBC menghapusnya dengan IV;
                  CFB dan OFB bekerja per unit $s$ bit; CTR memungkinkan
                  pemrosesan paralel.
            \item Confusion dicapai substitusi (S-box), diffusion dicapai
                  permutasi (P-box); Feistel membuat dekripsi tidak
                  memerlukan $f^{-1}$.
            \item DES memakai blok 64 bit, kunci efektif 56 bit, dan 16
                  putaran; ruang $2^{56}$ terlalu kecil, sehingga Double
                  DES tumbang oleh meet-in-the-middle dan Triple DES
                  memperpanjang kunci.
        \end{enumerate}
    \end{block}
    \vspace{0.2cm}
    \begin{exampleblock}{Pertemuan berikutnya}
        \textbf{Minggu 10 --- AES dan cipher blok lain}
    \end{exampleblock}
\end{frame}

\end{document}
